\chapter{La struttura della frase e il periodo}
\label{chap:struttura_frase}

La musica non è una successione casuale di accordi, ma un discorso organizzato. Le cadenze fungono da punteggiatura, definendo unità strutturali chiamate frasi e periodi.

\section{La frase musicale}
\label{sec:frase_musicale}

\begin{definizione}[Frase musicale]
    È l'unità minima di senso compiuto nel discorso tonale. Tipicamente composta da 4 misure (ma con varianti), termina sempre con una cadenza che ne delimita il confine.
\end{definizione}

\begin{regola}
    In un basso dato, ogni frase deve concludersi con una cadenza riconoscibile. Proprio come una frase verbale ha un soggetto, uno sviluppo e un punto finale, la frase musicale ha un inizio, un percorso armonico e una conclusione cadenzale.
\end{regola}

\section{Il periodo}
\label{sec:periodo}

\begin{definizione}[Il periodo]
    È una struttura di norma di 8 misure, formata dalla successione di due frasi in relazione tra loro:
    \begin{itemize}
        \item \textbf{Antecedente} (mis. 1-4): termina con una cadenza non conclusiva (sospesa o imperfetta). Rimanda a \ref{def:cadenza_sospesa} e \ref{def:cadenza_imperfetta}.
        \item \textbf{Conseguente} (mis. 5-8): termina con una cadenza conclusiva (perfetta). Rimanda a \ref{def:cadenza_perfetta}.
    \end{itemize}
\end{definizione}

La relazione tra antecedente e conseguente è analoga a quella tra domanda e risposta: la prima frase crea un'aspettativa, la seconda la soddisfa.

\begin{esempio}[Struttura del periodo in un basso dato]
    Indichiamo con \textit{markup} i punti cadenzali di un basso di 8 misure.
    \musica{es_cap07_01}
\end{esempio}

\section{Applicazione al basso dato}
\label{sec:applicazione_frase}

\begin{regola}[Pianificazione cadenzale]
    Prima di armonizzare un basso, individuare:
    \begin{itemize}
        \item \textbf{Metà (mis. 4 di 8):} Deve presentare una cadenza sospesa (V) o imperfetta (I non fondamentale o soprano non in ottava).
        \item \textbf{Finale (mis. 8):} Deve presentare una cadenza perfetta (V-I in stato fondamentale).
    \end{itemize}
\end{regola}

\begin{regola}[Scelta della cadenza di metà]
    Preferire la cadenza sospesa su V se il basso si ferma sulla dominante; preferire la cadenza imperfetta se il basso approda al I ma senza la forza conclusiva necessaria per il finale.
\end{regola}

\begin{esempio}[Armonizzazione di un periodo di 8 misure]
    Mostriamo come la scelta delle cadenze definisce la struttura formale del brano.
    \musica{es_cap07_02}
\end{esempio}
